% =============================================================
% Plain-Language Summary (optional lay summary; not currently
% \input from main.tex)
% =============================================================

Automated programs read medical research papers and extract facts
about how drugs, genes, and diseases relate to one another. These
facts feed cancer knowledge databases that doctors and researchers
rely on. Developers usually test such programs on standard benchmark
datasets to gauge accuracy. We asked whether scoring well on a
benchmark also means the program will reliably surface the right
facts for a real cancer knowledge base.

We trained 310 versions of a text-extraction program, systematically
varying four design choices, and tested each version against both a
standard benchmark and 162 expert-curated facts from the CIViC
oncology database. The two scores did not always agree: some
versions that looked equally good on the benchmark differed
substantially in how well they retrieved real cancer knowledge,
and the disagreement depended on which design choice was varied.
Because the evidence is mixed and statistical uncertainty is wide,
we recommend that developers check their programs against both a
benchmark and a small audit of expert-curated facts, rather than
relying on benchmark scores alone.
